\begin{textsample}{POS Dim 4 – 2010\_c4 – Score -1.00 – 2010\_c4/enade\_2010\_zoot\_22152f82a0.txt}  \label{ex:f4_pos_293}
CAPÍTULO 4 PERCEPÇÃO DA PROVA As análises feitas neste capítulo tratam das percepções dos estudantes, separados por grupos de Ingressantes e Concluintes, sobre a prova aplicada no ENADE/2010 na Área de Zootecnia.
Estas percepções foram mensuradas por meio de nove questões que avaliaram desde o grau de dificuldade da prova até o tempo gasto para concluí -la.
As percepções sobre a prova foram relacionadas com o desempenho dos alunos e com a Grande Região de funcionamento do curso.
O desempenho dos alunos foi classificado em dois níveis : quarto inferior e quarto superior.
Para tanto, esse desempenho foi colocado em ordem ascendente.
O percentil 25, P25, também conhecido como primeiro quartil, é a nota de desempenho que deixa um quarto ( 25,0 \% ) dos valores observados abaixo e três quartos acima.
A Figura 1 apresenta uma ilustração destes conceitos.
O quarto inferior de desempenho é composto pelas notas abaixo do primeiro quartil.
Já o percentil 75, P75, também conhecido como terceiro quartil, é o valor para o qual há três quartos ( 75,0 \% ) dos dados abaixo e um quarto acima dele.
O quarto superior de desempenho é composto pelas notas iguais ou acima do terceiro quartil.
O percentil 50, P50, também conhecido como mediana, é o valor que divide as notas em dois conjuntos de igual tamanho.
Vale ressaltar que percentis, quartis e medianas são pontos, e não obrigatoriamente pertencentes ao conjunto original de dados.
Ao passo que os quartos são subconjuntos dos dados originais.
A seguir, serão apresentados gráficos com resultados selecionados, relativos às nove questões avaliadas por grupos de estudantes.
As Tabelas no Anexo II apresentam os valores absolutos e a distribuição percentual das alternativas válidas das nove questões, segundo o mesmo recorte de desempenho dos alunos e Grande Região de funcionamento do curso. 4.1 GRAU DE DIFICULDADE DA PROVA 4.1.1 Componente de Formação Geral Na questão “ Qual o grau de dificuldade desta prova na parte de Formação Geral? ”, 19,6 \% do grupo de Ingressantes e 20,0 \% do grupo de Concluintes optaram pelas alternativas difícil ou muito difícil.
Entretanto, para mais da metade dos estudantes, tanto Ingressantes ( 63,4 \% ) quanto Concluintes ( 64,3 \% ), o Componente de Formação Geral da prova foi considerado com grau de dificuldade médio ( Gráfico 4.1 e Tabelas II.1 e II.10 no Anexo II ).
O percentual de Ingressantes que consideraram a prova difícil ou muito difícil foi mais elevado no quarto inferior ( 23,9 \% ) do que no quarto superior ( 16,5 \% ).
No grupo dos Concluintes, houve percentual de participantes maior no quarto superior ( 21,6 \% ), do que no quarto inferior ( 20,2 \% ).
Semelhantemente ao ocorrido para os Ingressantes como um todo, para os quartos de menor e maior desempenho deste grupo, a alternativa modal para esta pergunta foi médio, com respectivamente 62,0 \% e 64,3 \% dos respondentes.
Já entre os Concluintes, estas percentagens foram, respectivamente, 64,1 \% e 61,3 \%.
Considerando -se as alternativas escolhidas pelos estudantes de cada Grande Região, a região Sul foi aquela onde houve maior proporção, no grupo de Ingressantes, com 22,1 \%, enquanto a de menor incidência foi a região Sudeste, com 17,1 \% dos respondentes considerando a prova como difícil ou muito difícil.
Quanto ao grupo de Concluintes, que avaliaram o Componente de Formação Geral da prova da mesma forma, 25,1 \% da região Centro-Oeste tiveram, em maior medida, essa percepção, enquanto 16,7 \% da região Sudeste tiveram a mesma percepção, em menor frequência.
Nas Grandes Regiões, a proporção de Ingressantes que consideraram o Componente de Formação Geral da prova como sendo de grau de dificuldade médio esteve localizada entre 60,7 \% na região Sudeste e 68,1 \% na região Centro-Oeste.
A opção pela alternativa de grau médio por Grande Região que foi feita pelos Concluintes, atingiu a percentuais que estiveram entre 61,9 \% na região Sul e 69,6 \%, na região Norte. 4.1.2 Componente de Conhecimento Específico Ao responderem à questão “ Qual o grau de dificuldade desta prova na parte de Componente Específico? ”, 56,4 \% do grupo de Ingressantes e 26,5 \% do de Concluintes classificaram -na como difícil ou muito difícil.
O Componente de Conhecimento Específico da prova foi considerado com grau de dificuldade médio por 40,0 \% dos Ingressantes e por 63,5 \% dos Concluintes ( Gráfico 4.2 e Tabelas II.2 e II.11 no Anexo II ).
Considerando -se a percepção sobre o grau de dificuldade do Componente de Conhecimento Específico da prova de acordo com o desempenho dos estudantes, observa -se que 59,0 \% dos Ingressantes que obtiveram desempenho situado no quarto superior classificaram -na como difícil ou muito difícil.
Dos Ingressantes que obtiveram desempenho situado no quarto inferior, 51,5 \% optaram pela mesma alternativa.
Para os Concluintes, o percentual de estudantes que consideraram o Componente de Conhecimento Específico da prova difícil ou muito difícil foi um pouco maior entre os que obtiveram desempenho no quarto superior ( 26,9 \% ) do que os que obtiveram desempenho no quarto inferior ( 23,5 \% ).
Semelhantemente ao ocorrido para os Ingressantes como um todo, para os quartos de menor e maior desempenho deste grupo, a alternativa modal para esta pergunta foi difícil, com respectivamente 41,4 \% e 51,3 \% dos respondentes.
Já entre os Concluintes, a alternativa modal para a Questão 2 foi médio, com 63,4 \% do quarto inferior e 61,0 \% do quarto superior.
Entre as Grandes Regiões brasileiras, a proporção de Ingressantes que avaliaram o Componente de Conhecimento Específico da prova como difícil ou muito difícil variou de 51,6 \% ( Centro-Oeste ) a 60,9 \% ( Sudeste ).
O grupo de Concluintes avaliou esse componente como difícil ou muito difícil com menor frequência, com o mínimo alcançado na região Sudeste, 20,7 \%, e o máximo na região Norte, com 34,6 \% dos respondentes.
O percentual de Ingressantes que classificaram o grau de dificuldade como médio, no Componente de Conhecimento Específico da prova variou de 34,2 \% ( Sudeste ) a 43,4 \% ( Nordeste ).
Dentre os Concluintes, a percepção do mesmo grau de dificuldade aconteceu com maior frequência, variando de 54,6 \% ( Norte ) a 67,4 \% ( Sudeste ). 4.2 EXTENSÃO DA PROVA EM RELAÇÃO AO TEMPO TOTAL Indagados quanto à extensão da prova, em relação ao tempo total oferecido para a sua resolução, os estudantes apontaram, com maior incidência, a alternativa que considerava a extensão adequada, para ambos os grupos e todas as agregações consideradas ( Gráfico 4.3 e Tabelas II.3 e II.12 no Anexo II ).
No grupo de Ingressantes, a maioria, 64,1 \%, respondeu ser a extensão da prova adequada, ao passo que, no grupo de Concluintes, este valor foi menor, 48,1 \%.
Para 30,6 \% dos Ingressantes o exame foi longo ou muito longo.
Essa mesma consideração foi feita por 47,1 \% dos Concluintes.
Quanto aos extremos de desempenho, no grupo de Ingressantes, a maioria, tanto para o quarto inferior ( 63,4 \% ) quanto para o de melhor desempenho ( 65,3 \% ) respondeu ser a extensão da prova adequada.
Já entre os Concluintes, esses valores foram menores, situados em cada quarto em, respectivamente, 54,0 \% e 42,4 \%.
Considerando -se ainda os quartos extremos de desempenho, observa -se que a proporção de estudantes que consideraram a prova longa ou muito longa, em relação ao tempo total destinado à sua resolução, acompanhou a percepção de Ingressantes e Concluintes.
Para o grupo de Ingressantes, entre os estudantes no quarto de desempenho inferior, 31,1 \% declararam ter esta percepção, ao passo que entre aqueles no quarto de melhor desempenho, esta proporção foi 29,5 \%.
O mesmo pode ser constatado entre os Concluintes : 40,3 \% no quarto de menor desempenho e 54,3 \% no quarto de maior desempenho declararam ter tido esta percepção.
Na análise por Grandes Regiões, o percentual de Ingressantes que responderam ser a extensão da prova adequada variou de 61,0 \% no Norte a 70,3 \% no Centro-Oeste.
Já entre os Concluintes, estes valores estavam compreendidos entre 43,9 \% ( Nordeste ) e 53,3 \% ( Centro-Oeste ).
Aqueles Ingressantes que avaliaram a prova como longa ou muito longa em relação ao tempo total destinado à sua resolução variou de 22,5 \% na região Centro-Oeste a 33,6 \% na região Sudeste.
Os percentuais, quanto ao mesmo aspecto, entre os Concluintes, foram de 36,9 \% na região Norte, indo até 50,2 \% na região Nordeste. 4.3 COMPREENSÃO DOS ENUNCIADOS DAS QUESTÕES 4.3.1 Componente de Formação Geral Com relação aos enunciados das questões do Componente de Formação Geral, as opiniões foram positivas, com percentuais que foram dos 73,9 \% de Concluintes aos 79,7 \% de Ingressantes que consideraram claros e objetivos os enunciados de todas ou da maioria das questões ( Gráfico 4.4 e Tabelas II.4 e II.13 no Anexo II ).
Ao se comparar as percepções dos estudantes segundo o desempenho, observou -se que, entre os Ingressantes do quarto superior, a clareza e objetividade de todos ou da maioria dos enunciados das questões foram percebidas por 83,7 \%.
A proporção apurada para os Concluintes desse mesmo quarto foi 74,7 \%.
No quarto de desempenho inferior, tal avaliação foi emitida por 76,2 \% dos Ingressantes e por 69,0 \% dos Concluintes.
A partir dessas estatísticas, é possível afirmar que todos ou a maioria dos enunciados das questões do Componente de Formação Geral estavam claros e objetivos para a maioria dos respondentes, segundo os quartos de desempenho considerados, embora os que se situaram no quarto superior tenham declarado esta percepção com maior incidência.
Na análise regional, observa -se que as opiniões de que todos ou a maioria dos enunciados das questões do Componente de Formação Geral estavam claros e objetivos variaram, entre os Ingressantes, de 78,0 \% na região Norte a 81,0 \% na região Sudeste.
Entre os Concluintes os valores correspondentes variaram de 71,6 \% ( na região Norte ) a 76,1 \% ( na região Sul ). 4.3.2 Componente de Conhecimento Específico Com relação aos enunciados das questões do Componente de Conhecimento Específico da prova, para os Concluintes da Área de Zootecnia, a clareza e a objetividade foram percebidas em todos ou na maioria destes por 78,9 \%.
Para o grupo dos Ingressantes, foram percebidas por 75,2 \%.
A proporção de estudantes que considerou os enunciados das questões claros e objetivos foi mais elevada nos quartos superiores : 77,7 \% dos Ingressantes e 80,7 \% dos Concluintes, se comparados aos 70,5 \% e aos 76,1 \% dos mesmos grupos, respectivamente, que estiveram no quarto inferior de desempenho ( Gráfico 4.5 e Tabelas II.5 e II.14 no Anexo II ).
Conhecimento Específico da prova variou entre 72,4 \% ( Norte ) e 80,8 \% ( Centro-Oeste ).
Entre os Ingressantes, os índices variaram entre 70,7 \% na região Norte e 79,1 \%, também na região Centro-Oeste. 4.4 SUFICIÊNCIA DAS INFORMAÇÕES/INSTRUÇÕES FORNECIDAS Ao avaliarem as informações/instruções fornecidas para a resolução das questões, 81,5 \% do grupo dos Ingressantes e 88,1 \% do grupo dos Concluintes respondentes de todo o Brasil afirmaram que estas eram até excessivas ou suficientes em todas ou na maioria das questões ( Gráfico 4.6 e Tabelas II.6 e II.15 no Anexo II ).
Levando -se em conta o desempenho dos participantes, notam -se diferenças entre as opiniões de estudantes, Ingressantes e Concluintes, dos quartos inferior e superior de desempenho, como mostra o Gráfico 4.6.
O percentual de participantes que avaliaram as informações/instruções como excessivas ou suficientes em todas ou na maioria das questões foi mais elevado nos quartos superiores, correspondendo a 85,0 \% dos Ingressantes e a 91,9 \% dos Concluintes, percentuais superiores à média nacional para cada um dos dois grupos.
Já nos quartos inferiores, a suficiência das informações/instruções declarada como excessiva em todas ou na maioria das questões foi percebida por 79,1 \% dos Ingressantes e por 81,2 \% dos Concluintes.
Quanto à distribuição de respondentes nas Grandes Regiões, observa -se que as proporções de estudantes que perceberam como excessivas ou suficientes em todas ou na maioria das questões as informações/instruções fornecidas para a resolução das mesmas superaram as médias nacionais, em maior medida, entre os Ingressantes da região Centro-Oeste ( 85,6 \% ).
Entre os Concluintes, a superação da média nacional foi mais expressiva na região Sul ( 93,0 \% ). 4.5 DIFICULDADE ENCONTRADA AO RESPONDER À PROVA Perguntados sobre as dificuldades com as quais se depararam ao responder à prova, 62,7 \% do grupo de Ingressantes e 10,1 \% do grupo de Concluintes apontaram o desconhecimento do conteúdo.
Concluintes, por sua vez, apontaram a forma diferente de abordagem do conteúdo como dificuldade, mencionada por 43,2 \% deste grupo.
Esta alternativa foi mencionada por uma parcela menor dos Ingressantes, 18,0 \%.
A falta de motivação para fazer a prova foi a dificuldade para 21,8 \% dos Concluintes e para 9,8 \% dos Ingressantes.
Considerando -se os respondentes em todo o Brasil, 7,4 \% dos Ingressantes e 17,8 \% dos Concluintes afirmaram que não tiveram qualquer tipo de dificuldade para responder à prova ( Tabelas II.7 e II.16 no Anexo II ).
Com relação aos quartos de desempenho, o desconhecimento do conteúdo foi a dificuldade modal entre os Ingressantes, apontada por 63,8 \% daqueles no quarto superior e por 54,8 \% dos Ingressantes do quarto inferior.
Entre os Concluintes, os percentuais de estudantes com a mesma percepção foi inferior ao encontrado entre os Ingressantes, 8,4 \% no quarto superior e 9,7 \% no quarto inferior.
A alternativa modal pelo grupo de Concluintes, quando desagregada pelo desempenho, foi que a dificuldade encontrada era a forma diferente de abordagem do conteúdo : 39,7 \% no quarto inferior e 44,9 \% do quarto superior assim o responderam.
Na análise por Grandes Regiões, o percentual de Ingressantes que apontaram o desconhecimento do conteúdo como dificuldade ao responder à prova variou de 55,2 \%, na região Norte a 69,1 \%, na Sudeste.
Entre os Concluintes, tal percentual foi mais baixo, variando de 8,2 \% na região Sudeste a 13,6 \%, na Norte.
A forma diferente de abordagem do conteúdo foi citada, entre os Concluintes, por percentuais que variaram entre 36,8 \% ( região Sudeste ) e 52,8 \% ( Norte ).
Em contrapartida, os percentuais relacionados a essa opção entre os Ingressantes foi mais baixo, variando de 12,8 \% ( região Sudeste ) a 29,8 \% ( região Norte ).
O percentual de Ingressantes que citou a falta de motivação como dificuldade superou a média nacional, em maior medida, na região Nordeste ( 11,1 \% ).
Entre os Concluintes esse percentual variou entre 13,6 \% ( região Norte ) e 28,8 \% ( região Sudeste ).
Houve percentual acima da média nacional dos Ingressantes nas regiões Nordeste ( 8,8 \% ) e Centro-Oeste ( 9,2 \% ), com declarações de que não tiveram qualquer dificuldade para responder à prova.
Entre os Concluintes, esta situação ocorreu, em maior escala, na região Sul, com 22,8 \% de respondentes.
O Gráfico 4.7 apresenta os percentuais de estudantes, Ingressantes e Concluintes, que apontaram o desconhecimento do conteúdo como dificuldade percebida ao responder à prova. 4.6 CONTEÚDOS DAS QUESTÕES OBJETIVAS DA PROVA Ao analisarem os conteúdos das questões objetivas da prova, a maior parte dos Ingressantes, 81,1 \% do total nacional, afirmou que não estudou ainda a maioria desses conteúdos, enquanto entre os Concluintes, apenas 2,0 \% expressaram a mesma opinião.
Por outro lado, a maioria dos Concluintes ( 80,1 \% ) afirmou ter estudado e aprendido muitos ou todos os conteúdos.
Para os Ingressantes, o percentual correspondente a essas alternativas foi 10,1 \% ( Gráfico 4.8 e Tabelas II.8 e II.17 no Anexo II ).
Considerando -se separadamente as opiniões de estudantes dos quartos de desempenho inferior e superior, observa -se que, para os Ingressantes, o percentual de participantes que afirmaram que não estudaram ainda a maioria desses conteúdos, foi um pouco mais elevado no quarto superior ( 79,8 \% ) do que no inferior ( 77,9 \% ).
Para os Concluintes do quarto inferior, 4,9 \% ofereceram essa resposta, sendo 1,0 \% os do quarto superior.
Tendo em conta os quartos superiores para os dois grupos, 91,4 \% dos Concluintes e 13,3 \% dos Ingressantes afirmaram ter estudado e aprendido muitos ou todos os conteúdos.
Na análise por Grande Região, a proporção de Ingressantes que afirmou que não estudou ainda a maioria desses conteúdos superou a média nacional nas regiões Sudeste ( 84,5 \% ) e Norte ( 85,1 \% ).
Já entre os Concluintes, as regiões Norte e Centro-Oeste apresentaram percentuais maiores do que a média nacional para esta alternativa, ambas com 3,6 \%.
Em todas as regiões, percentuais acima da metade dos Concluintes afirmaram ter estudado e aprendido muitos ou todos os conteúdos, variando entre 71,4 \% na região Norte e 86,5 \% na região Sul. 4.7 TEMPO GASTO PARA CONCLUIR A PROVA Ao responderem sobre o tempo de conclusão da prova, mais da metade dos estudantes afirmou ter gasto entre duas e quatro horas.
Considerando -se separadamente as percepções do grupo de Ingressantes e do grupo de Concluintes, os percentuais, em termos nacionais, foram 63,0 \% e 71,3 \%, respectivamente ( Gráfico 4.9 e Tabelas II.9 e II.18 no Anexo II ).
Uma vez sendo analisadas as alternativas oferecidas pelos estudantes que se situaram nos quartos superior e inferior, observa -se que uma maior proporção de Ingressantes no quarto superior declarou ter gasto entre duas e quatro horas para concluir a prova quando comparados com os do quarto inferior, respectivamente, 76,8 \% e 46,6 \%.
A mesma situação foi verificada no grupo dos Concluintes, no qual 81,9 \% daqueles no quarto superior e 56,4 \% daqueles no quarto inferior declaram ter gasto entre duas e quatro horas para concluir a prova.
Nas cinco Grandes Regiões brasileiras, os Ingressantes que utilizaram entre duas e quatro horas para finalizar a prova superaram o percentual nacional nas regiões Norte ( 64,3 \% ), Nordeste ( 65,1 \% ) e Centro-Oeste, ( 67,3 \% ).
Entre os Concluintes, o mesmo foi verificado entre os estudantes das regiões Sul ( 76,3 \% ) e Centro-Oeste ( 73,5 \% ).
A região Norte foi a que apresentou o menor percentual de Concluintes que dispensaram entre duas e quatro horas para concluir a prova : 61,3 \%, como mostra o Gráfico 4.9.

% matched lemmas: 
\end{textsample}
